% sec8_4.tex — Section 8.4: Multivariate Regressions

In Section~1.5 and also in Example~7.2 of Chapter~7, we pointed out that the OLS
estimator of the linear regression model is numerically the same as the ML estimator of the same model with normal errors. We have derived the GMM estimators
for the single-equation models with endogenous regressors in Chapter~3 and for
their multiple-equation versions in Chapter~4, but we have not indicated what the
ML estimators for those models are. In this section and the next, we derive those
ML counterparts. This section examines the ML estimation of the multivariate
regression model, which is the simplest multiple-equation model.

\subsection*{The Multivariate Regression Model Restated}

The multivariate regression model is described by Assumptions~4.1--4.5 and~4.7,
with $\mathbf{z}_{tm} = \mathbf{x}_t$ for all $m = 1, 2, \ldots, M$. To change the notation of Chapter~4,
the system of $M$ equations can be written as
\begin{equation}
  y_{tm} = \underset{(1 \times K)}{\mathbf{x}_t'} \;
  \underset{(K \times 1)}{\boldsymbol{\pi}_{0m}}
  + v_{tm}
  \quad (m = 1, 2, \ldots, M;\; t = 1, 2, \ldots, n).
  \label{eq:8.4.1} \tag{8.4.1}
\end{equation}
Here, we maintain the convention of designating the true parameter vector by subscript~0. If we define
\begin{equation}
  \underset{(M \times 1)}{\mathbf{y}_t}
  = \begin{bmatrix} y_{t1} \\ y_{t2} \\ \vdots \\ y_{tM} \end{bmatrix}, \quad
  \underset{(M \times 1)}{\mathbf{v}_t}
  = \begin{bmatrix} v_{t1} \\ v_{t2} \\ \vdots \\ v_{tM} \end{bmatrix}, \quad
  \underset{(K \times M)}{\boldsymbol{\Pi}_0}
  = \begin{bmatrix}
    \boldsymbol{\pi}_{01} & \boldsymbol{\pi}_{02} & \cdots & \boldsymbol{\pi}_{0M}
  \end{bmatrix},
  \label{eq:8.4.2} \tag{8.4.2}
\end{equation}
then the $M$-equation system can be written compactly as
\begin{equation}
  \mathbf{y}_t = \boldsymbol{\Pi}_0'\mathbf{x}_t + \mathbf{v}_t
  \quad (t = 1, 2, \ldots, n).
  \label{eq:8.4.3} \tag{8.4.3}
\end{equation}
The model has the following features.
\begin{itemize}
\item $\{y_t, \mathbf{x}_t\}$ is ergodic stationary (Assumption~4.2).
\item The common regressors $\mathbf{x}_t$ are predetermined in that
  $\E(\mathbf{x}_t \cdot v_{tm}) = \mathbf{0}$ for all~$m$ (Assumption~4.3).
\item The rank condition for identification (Assumption~4.4) reduces to the familiar
  condition that $\E(\mathbf{x}_t\mathbf{x}_t')$ be nonsingular.
\item The error vector is conditionally homoskedastic in that
  $\E(\mathbf{v}_t\mathbf{v}_t' \mid \mathbf{x}_t) = \boldsymbol{\Omega}_0$, and
  $\boldsymbol{\Omega}_0$ is positive definite (Assumption~4.7).
\end{itemize}

The GMM estimator of $\boldsymbol{\pi}_{0m}$ is the OLS estimator:
\begin{equation}
  \hat{\boldsymbol{\pi}}_m
  = \biggl(\frac{1}{n}\sum_{t=1}^{n} \mathbf{x}_t\mathbf{x}_t'\biggr)^{-1}
    \biggl(\frac{1}{n}\sum_{t=1}^{n} \mathbf{x}_t \cdot y_{tm}\biggr).
  \label{eq:8.4.4} \tag{8.4.4}
\end{equation}
So the GMM (OLS) estimate of $\boldsymbol{\Pi}_0$ is given by
\begin{equation}
  \underset{(K \times M)}{\hat{\boldsymbol{\Pi}}}
  = \biggl(\frac{1}{n}\sum_{t=1}^{n} \mathbf{x}_t\mathbf{x}_t'\biggr)^{-1}
    \underset{(K \times K)}{\phantom{x}}
    \biggl(\frac{1}{n}\sum_{t=1}^{n} \mathbf{x}_t\mathbf{y}_t'\biggr).
    \underset{(K \times M)}{\phantom{x}}
  \label{eq:8.4.5} \tag{8.4.5}
\end{equation}

\subsection*{The Likelihood Function}

The model is not yet specific enough to lend itself to ML estimation. We make the
supplementary assumptions that (1) $\mathbf{v}_t \mid \mathbf{x}_t \sim N(\mathbf{0}, \boldsymbol{\Omega}_0)$, and
(2) $\{y_t, \mathbf{x}_t\}$ is i.i.d.
(This assumption strengthens Assumptions~4.2 and~4.5.)

The multivariate regression model with the normality assumption~(1) implies
that $\mathbf{y}_t \mid \mathbf{x}_t \sim N(\boldsymbol{\Pi}_0'\mathbf{x}_t, \boldsymbol{\Omega}_0)$.
The density of this multivariate normal distribution is
\begin{equation}
  (2\pi)^{-M/2}\,|\boldsymbol{\Omega}_0|^{-1/2}
  \exp\biggl[-\frac{1}{2}(\mathbf{y}_t - \boldsymbol{\Pi}_0'\mathbf{x}_t)'
    \boldsymbol{\Omega}_0^{-1}
    (\mathbf{y}_t - \boldsymbol{\Pi}_0'\mathbf{x}_t)\biggr].
  \label{eq:8.4.6} \tag{8.4.6}
\end{equation}
Replacing the true parameter values $(\boldsymbol{\Pi}_0, \boldsymbol{\Omega}_0)$ by their hypothetical values
$(\boldsymbol{\Pi}, \boldsymbol{\Omega})$ and taking logs, we obtain the log conditional likelihood for observation~$t$:
\begin{multline}
  \log f(\mathbf{y}_t \mid \mathbf{x}_t;\, \boldsymbol{\Pi}, \boldsymbol{\Omega}) \\
  = -\frac{M}{2}\log(2\pi)
    + \frac{1}{2}\log(|\boldsymbol{\Omega}^{-1}|)
    - \frac{1}{2}(\mathbf{y}_t - \boldsymbol{\Pi}'\mathbf{x}_t)'
      \boldsymbol{\Omega}^{-1}
      (\mathbf{y}_t - \boldsymbol{\Pi}'\mathbf{x}_t),
  \label{eq:8.4.7} \tag{8.4.7}
\end{multline}
where use has been made of the fact that $-\log(|\boldsymbol{\Omega}|) = \log(|\boldsymbol{\Omega}^{-1}|)$.

For a random sample, $1/n$ times the log conditional likelihood of the sample
is the average over~$t$ of the log conditional likelihood for observation~$t$. So the
objective function $Q_n(\boldsymbol{\Pi}, \boldsymbol{\Omega})$ is
\begin{equation}
  Q_n(\boldsymbol{\Pi}, \boldsymbol{\Omega})
  = -\frac{M}{2}\log(2\pi)
    + \frac{1}{2}\log(|\boldsymbol{\Omega}^{-1}|)
    - \frac{1}{2n}\sum_{t=1}^{n}
      (\mathbf{y}_t - \boldsymbol{\Pi}'\mathbf{x}_t)'
      \boldsymbol{\Omega}^{-1}
      (\mathbf{y}_t - \boldsymbol{\Pi}'\mathbf{x}_t).
  \label{eq:8.4.8} \tag{8.4.8}
\end{equation}
It is left as Review Question~2 to show that the last term can be written as
\begin{equation}
  \frac{1}{2n}\sum_{t=1}^{n}
    (\mathbf{y}_t - \boldsymbol{\Pi}'\mathbf{x}_t)'
    \boldsymbol{\Omega}^{-1}
    (\mathbf{y}_t - \boldsymbol{\Pi}'\mathbf{x}_t)
  = \frac{1}{2}\,\tr[\boldsymbol{\Omega}^{-1}\hat{\boldsymbol{\Omega}}(\boldsymbol{\Pi})],
  \label{eq:8.4.9} \tag{8.4.9}
\end{equation}
where $\hat{\boldsymbol{\Omega}}(\boldsymbol{\Pi})$ is defined as
\begin{equation}
  \underset{(M \times M)}{\hat{\boldsymbol{\Omega}}(\boldsymbol{\Pi})}
  \equiv \frac{1}{n}\sum_{t=1}^{n}
    (\mathbf{y}_t - \boldsymbol{\Pi}'\mathbf{x}_t)
    (\mathbf{y}_t - \boldsymbol{\Pi}'\mathbf{x}_t)'.
  \label{eq:8.4.10} \tag{8.4.10}
\end{equation}
So the objective function can be rewritten as
\begin{equation}
  Q_n(\boldsymbol{\Pi}, \boldsymbol{\Omega})
  = -\frac{M}{2}\log(2\pi)
    + \frac{1}{2}\log(|\boldsymbol{\Omega}^{-1}|)
    - \frac{1}{2}\,\tr[\boldsymbol{\Omega}^{-1}\hat{\boldsymbol{\Omega}}(\boldsymbol{\Pi})].
  \label{eq:8.4.11} \tag{8.4.11}
\end{equation}
The ML estimate of $(\boldsymbol{\Pi}_0, \boldsymbol{\Omega}_0)$ is the $(\boldsymbol{\Pi}, \boldsymbol{\Omega})$ that maximizes this objective function.

The parameter space for $(\boldsymbol{\Pi}_0, \boldsymbol{\Omega}_0)$ is
\[
  \{(\boldsymbol{\Pi}, \boldsymbol{\Omega}) \mid \boldsymbol{\Omega} \text{ is symmetric and positive definite}\}.
\]

\subsection*{Maximizing the Likelihood Function}

We now prove that the solution to the maximization problem is numerically the
same as the GMM estimator. That is, the ML estimator of $\boldsymbol{\Pi}_0$ is given by the
GMM/OLS estimator $\hat{\boldsymbol{\Pi}}$ in \eqref{eq:8.4.5} and the ML estimate of $\boldsymbol{\Omega}_0$ is given by
\begin{equation}
  \hat{\boldsymbol{\Omega}}
  = \frac{1}{n}\sum_{t=1}^{n}
    \hat{\mathbf{v}}_t \hat{\mathbf{v}}_t'
  = \hat{\boldsymbol{\Omega}}(\hat{\boldsymbol{\Pi}}),
  \label{eq:8.4.12} \tag{8.4.12}
\end{equation}
where $\hat{\mathbf{v}}_t \equiv \mathbf{y}_t - \hat{\boldsymbol{\Pi}}'\mathbf{x}_t$.

It is left as Analytical Exercise~2 to show that, under the assumptions of the
multivariate regression model stated above, $\hat{\boldsymbol{\Omega}}(\boldsymbol{\Pi})$ is positive definite with probability one for any given $\boldsymbol{\Pi}$ for sufficiently large~$n$. So we can assume that
$\hat{\boldsymbol{\Omega}}(\boldsymbol{\Pi})$ is positive definite, not just positive semidefinite.
The proof that the GMM estimator
is numerically the same as the ML estimator is based on the following two-step
maximization of the objective function.

\medskip
\noindent\textit{Step 1:}\quad
The first step is to maximize $Q_n(\boldsymbol{\Pi}, \boldsymbol{\Omega})$ with respect to $\boldsymbol{\Omega}$, taking $\boldsymbol{\Pi}$ as
given. For this purpose, the following fact from matrix algebra is useful.

\begin{quote}
\textit{An Inequality Involving Trace and Determinant:}\footnote{This result is implied by Lemma~A.6 of Johansen (1995). The uniqueness of the minimizer is due to the
linearity of the trace operator and the strict concavity of $\log(|\mathbf{A}|)$ noted in Magnus and Neudecker (1988, p.~222).}\quad
Let $\mathbf{A}$ and $\mathbf{B}$
be two symmetric and positive definite matrices of the same size.
Then the function
\[
  f(\mathbf{A}) \equiv \log(|\mathbf{A}|) - \tr(\mathbf{A}\mathbf{B})
\]
is maximized uniquely by $\mathbf{A} = \mathbf{B}^{-1}$.
\end{quote}

This result, with $\mathbf{A} = \boldsymbol{\Omega}^{-1}$ and
$\mathbf{B} = \hat{\boldsymbol{\Omega}}(\boldsymbol{\Pi})$, immediately implies that the
objective function \eqref{eq:8.4.11} is maximized uniquely by
$\boldsymbol{\Omega} = \hat{\boldsymbol{\Omega}}(\boldsymbol{\Pi})$, given $\boldsymbol{\Pi}$.
Substituting this $\boldsymbol{\Omega}$ into \eqref{eq:8.4.11} gives ($1/n$ times) the \textbf{concentrated
log likelihood function} (concentrated with respect to $\boldsymbol{\Omega}$):
\begin{align}
  Q_n^*(\boldsymbol{\Pi})
  &\equiv Q_n(\boldsymbol{\Pi},\, \hat{\boldsymbol{\Omega}}(\boldsymbol{\Pi})) \notag \\
  &= -\frac{M}{2}\log(2\pi)
     + \frac{1}{2}\log(|\hat{\boldsymbol{\Omega}}(\boldsymbol{\Pi})^{-1}|)
     - \frac{1}{2}\,\tr[\hat{\boldsymbol{\Omega}}(\boldsymbol{\Pi})^{-1}
       \hat{\boldsymbol{\Omega}}(\boldsymbol{\Pi})] \notag \\
  &= -\frac{M}{2}\log(2\pi)
     - \frac{1}{2}\log(|\hat{\boldsymbol{\Omega}}(\boldsymbol{\Pi})|)
     - \frac{M}{2}.
  \label{eq:8.4.13} \tag{8.4.13}
\end{align}

\medskip
\noindent\textit{Step 2:}\quad
Looking at the concentrated log likelihood function $Q_n^*(\boldsymbol{\Pi})$, we see that
the ML estimator of $\boldsymbol{\Pi}_0$ should minimize
\begin{equation}
  |\hat{\boldsymbol{\Omega}}(\boldsymbol{\Pi})|
  = \biggl|\frac{1}{n}\sum_{t=1}^{n}
    (\mathbf{y}_t - \boldsymbol{\Pi}'\mathbf{x}_t)
    (\mathbf{y}_t - \boldsymbol{\Pi}'\mathbf{x}_t)'\biggr|.
  \label{eq:8.4.14} \tag{8.4.14}
\end{equation}
It is left as Analytical Exercise~1 to show that this is minimized by the
OLS estimator $\hat{\boldsymbol{\Pi}}$ given in \eqref{eq:8.4.5}. Finally, setting
$\boldsymbol{\Pi} = \hat{\boldsymbol{\Pi}}$ in \eqref{eq:8.4.10}
gives \eqref{eq:8.4.12}.

\subsection*{Consistency and Asymptotic Normality}

We have shown in Chapter~4, without the normality assumption (that $\varepsilon_t \mid \mathbf{x}_t$ is
normal) or the i.i.d.\ assumption for $\{y_t, \mathbf{x}_t\}$, that the GMM/OLS estimator
$\hat{\boldsymbol{\Pi}}$ in \eqref{eq:8.4.5} is consistent and asymptotically normal and that $\hat{\boldsymbol{\Omega}}$ in
\eqref{eq:8.4.12} is consistent.
It then follows, trivially, that the ML estimator of $\boldsymbol{\Pi}_0$ is consistent and asymptotically normal and the ML estimator of $\boldsymbol{\Omega}_0$ is consistent. It is worth emphasizing
that the ML estimator of $\boldsymbol{\Pi}_0$ based on the likelihood function that assumes normality is consistent and asymptotically normal even if the error term is not normally
distributed.

We will not be concerned with the asymptotic normality of the ML estimator
of $\boldsymbol{\Omega}_0$, $\hat{\boldsymbol{\Omega}}(\hat{\boldsymbol{\Pi}})$. One way to demonstrate the asymptotic normality would be to verify
the conditions of the relevant asymptotic normality theorem of the previous chapter.

\bigskip
\noindent\rule{\textwidth}{0.4pt}
\textbf{QUESTION FOR REVIEW}
\medskip

\begin{enumerate}
\item Prove \eqref{eq:8.4.9}. \textbf{Hint:} Use the following properties of the trace operator.
  (1) $\tr(x) = x$ if $x$ is a scalar,
  (2) $\tr(\mathbf{A} + \mathbf{B}) = \tr(\mathbf{A}) + \tr(\mathbf{B})$, and
  (3) $\tr(\mathbf{A}\mathbf{B}) = \tr(\mathbf{B}\mathbf{A})$ provided $\mathbf{A}\mathbf{B}$ and $\mathbf{B}\mathbf{A}$
  are well defined. The answer is
  \begin{align*}
    \frac{1}{n}\sum_{t=1}^{n}
      &(\mathbf{y}_t - \boldsymbol{\Pi}'\mathbf{x}_t)'
      \boldsymbol{\Omega}^{-1}
      (\mathbf{y}_t - \boldsymbol{\Pi}'\mathbf{x}_t) \\
    &= \tr\biggl[\frac{1}{n}\sum_{t=1}^{n}
      (\mathbf{y}_t - \boldsymbol{\Pi}'\mathbf{x}_t)'
      \boldsymbol{\Omega}^{-1}
      (\mathbf{y}_t - \boldsymbol{\Pi}'\mathbf{x}_t)\biggr] \\
    &\phantom{= \tr\biggl[} \vdots
  \end{align*}
\end{enumerate}
